\documentclass[12pt, letterpaper]{article}

\setlength{\topmargin}{-1.75cm} \setlength{\textheight}{22.5cm}
\setlength{\oddsidemargin}{0.25cm}
\setlength{\evensidemargin}{0.25cm} \setlength{\textwidth}{16.2cm}

% symbols 
\usepackage{fdsymbol}
\usepackage{stmaryrd}
\usepackage{pifont} % check/cross marks 
\newcommand{\cmark}{\ding{51}}%
\newcommand{\xmark}{\ding{55}}%
%\usepackage{amssymb}
\usepackage[ruled,linesnumbered]{algorithm2e}
\usepackage{graphicx}
\usepackage[linguistics]{forest}
% for automata 
\usepackage{tikz} 
\usetikzlibrary{automata,positioning}


\usepackage{palatino, url, multicol} % for multiple columns

\usepackage{pictex}
%% in the .pictex output of xfig, there is command \colo
%% however the old version of pictex may not define this
%% so we define color here as empty
\def \color#1]#2{}

\usepackage{fancyhdr}
\pagestyle{fancy}

\newcommand{\hwnumber} {
    4
}

\newcommand{\hwtitle} {
  NFA. Context Free Grammar.  
}

\newcommand{\duedate}{
    11:59pm Wed Nov 5.
}

\date{}

\begin{document}

\newcommand{\hide}[1]{}
\newcommand{\otherquestions}[1]{}
\newcommand{\set}[1]{\{#1\}}
\newcommand{\pg}[1]{{\tt #1}}
\newtheorem{definition}{Definition}
\newcommand{\emptyclause}{\Box}
\newcommand{\keyword}[1]{{\bf #1}}

\newcommand{\ee}[1] {
  \begin{enumerate}
    #1
  \end{enumerate}
}
%% \noindent test \hfill test

\newcommand{\ie}[1] {
  \begin{enumerate}
    #1
  \end{enumerate}
}

\newcommand{\kleene}{\wedge}

\title{{\bf Homework \hwnumber.} \hwtitle}

\maketitle

\input{hwHeader}

 \noindent {\bf Submit PDF file (and Latex if you use it) of your solution to Canvas by
  \duedate.}

\ee{

    \item  You must use Latex.
    % For Latex, we will use the same policy as in HW2 and HW3. In short, everyone has a chance to practice Latex and get extra credit depending on the quality of the Latex source.  

	\item (25 - basics) (NFA) Consider the NFA $M = (\Sigma, K, s, F, \Delta)$ where $\Sigma = \{a, b, c\}$, $K = \{p_1, p_2, p_3\}$, $s = p_1$, $F=\{p3\}$  and $\Delta$ is a transition relation.   
	\ee{
		\item What is the alphabet of $M$? 
		\item What are the states of  $M$? 
		\item What is the start state of $M$? 
		\item Is $p_2$ a final state of $M$? 
		\item Give an example string over alphabet of $M$. 
		\item Give an example of a configuration of $M$ with state $p_1$ in the configuration. 
		\item Give an example of $\Delta$ with at least 1 element and at most 3 elements (recall that $\Delta$ is simply a set). 
	}

  \item (10 - basics) (Foundation) Study the {\em pumping lemma}. 
     \ee{
		\item List all concepts used in the lemma. Each concept should be in the format of {\em concept name} and {\em arguments} of the concept. 
		\item List all logical symbols used in the lemma. 
	}   % Question: definitions of concepts, write statement / (Q5 of HW2) for pumping lemma, list all concepts there in the form: concept name and their arguments. List all logical connectives there.  

   % 80% 
  \item (10 - medium) (Ideas underlying Pumping Lemma) Given an NFA $M$ and a computation over it:  $(p_1, abc)    \vdash_M  (p_2, bc) \vdash_M (p_2, c) \vdash_M (p_3, e)$ where $p_1$ is the start state and $p_3$ a final state. (Hint. Not only recall what {\em accept} intuitively means, but also its formal definition. Note also the repetition of state $p_2$ in the computation.)
	\ee{ 
		\item is string {\em abc} accepted by $M$? 
		\item is string {\em ac} accepted by $M$?   
		\item is string $ab^2c$ accepted by $M$?   
		\item Explain why for any non-negative number $n$,  $ab^nc$ is accepted by $M$? 
	}

  \item (5 -  advanced) (Comprehensive). Using pumping lemma, prove  the language $\{0^n1^n | n \ge 0\}$ over alphabet $\Sigma = \{0, 1\}$ is not regular. 

 \item (35 - basics) (Context free grammar). Use the problem decomposition method to 
	\ee{
		\item write  a precise definition AND a context free grammar for the language of all strings whose number of 1's is one less than its number of 0's.  % palindromes over alphabet $\{0, 1\}$.  Write also the English definition before your context free grammar. 
	} 
	and based on the grammar,
	\ee{
		\item Write the right most derivation for string  $100$. 
		\item Draw a parse tree for 100. 
	}
   
	% 90% 
	\item (5 - advanced) (Context free grammar). Following the method discussed in class, write context free grammar for the language (over alphabet $\{0, 1\}$) consisting of strings that have equal number 0's and 1's. If you follow ideas and methods other than the ones discussed in class, you will get at most 90\% for this quesiton. 

	% 80% 
	\item (10 - medium) (Foundation) Give a precise definition of {\em context free language}. 
}
\end{document}

